\scanpage{11}
% margin note written sideways
\margin{Catégorie ou \unsure{isotopique|multigroupe|topologique} des tubes \unsure{orientés|orientables|orientées}}
\(\left\{\vcenter{\hbox{\begin{minipage}{.89\linewidth}
\textcircled{a} $(\Sigma,\omega)$ $1$-complexe combinatoire \inserted{(fini)}, circularisé ``ordinaire'' (ss pts isolés ni boucles isolés)

\textcircled{c} $\Pi_d$ $\mathbb Z$-groupoïde \deleted{\unsure{(à $\pi_0(\Pi_b)$ fini)|(à $\pi_0(\Pi_d)$ fini)}} $1$-spécial

\textcircled{b} $\Pi_b$ groupoïde $1$-spécial (i.e. $\pi_0(\Pi_b)$ fini, les groupes fond. isom à $\mathbb Z$)


\end{minipage}}}\right.\)

\textcircled{a$'$} \(\left\{\vcenter{\hbox{\begin{minipage}{.89\linewidth}
$\deleted{\unsure{\Pi}}$ $\partial_t(\Sigma,\omega)$ $1$-complexe combinatoire \deleted{\unsure{circul.}} \uline{orienté} dit ``des co-circuits'' de $\Sigma$.
\[
 \Pi(\partial_t(\Sigma,\omega))\qquad \mathbb Z\text{-groupoïde }1\text{-spécial}
 \quad\xrightarrow{\varphi_0}\Pi(\Sigma)
\]


\end{minipage}}}\right.\)

\textcircled{b$'$}
\(\left\{\vcenter{\hbox{\begin{minipage}{.89\linewidth}

\[
 \overset{\deleted{\unsure{\sim}}}{\Pi}_b=\coprod_{i\in B}
 \underset{\text{comp. connexes}}{\underline{\Pi_b(i)}}
 \qquad \Lambda(i)=\pi_1\Pi_b(i)
\]
\hfill (groupe isom. à $\mathbb Z$)
\[
 \widetilde\Pi_b(i)=\deleted{\unsure{\mathbb Z}}
 \coprod_{\alpha\in\mathrm{Isom}_{\rm gr}(\mathbb Z,\Lambda_{\deleted{\unsure{b}}}(i))}
 \Pi_b(i)
\]
$\mathbb Z$-groupoïde \deleted{\unsure{connexe|grossier}} \inserted{à deux} composantes connexes
\diagram{\begin{tikzcd}[ampersand replacement=\&]
\widetilde\Pi_b \arrow[d,"\varphi_b"] \arrow[r,phantom,"="] \&
\displaystyle\coprod_{i\in B}\widetilde\Pi_b(i) \& \mathbb Z\text{-groupoïde }1\text{-spécial}\\
\Pi_b \& \&
\end{tikzcd}}


\end{minipage}}}\right.\)

\textcircled{c$'$}
$\mathrm{Cat}\,\pi_0(\Pi_d))\xleftarrow{\varphi_c}\Pi_d$

\[
 T(\Sigma,\omega,\Pi_d,\Pi_b)=(C,\underset{\nearrow\,\mathbb Z\text{-groupoïde }1\text{-spécial}}{D},\varphi)
\]
\[
\begin{aligned}
D&=\Pi(\partial_t(\Sigma,\omega))\amalg\widetilde\Pi_b\amalg\Pi_d\\
C&=\Pi(\Sigma)\amalg\Pi_b\amalg\mathrm{Cat}\,(\pi_0(\Pi_d))\\
\varphi&=\varphi_0\amalg\varphi_b\amalg\varphi_d
\end{aligned}
\]
\margin{\uline{NB} $\varphi$ est foncteur bordant}

\noindent\rule{\linewidth}{.4pt}

\textcircled{d} $C'\unsure{,}$ groupoïde

\textcircled{e} $\varphi':D\longrightarrow C'$ \qquad $(\varphi'=(\varphi'_0,\varphi'_b,\varphi'_d))$

\hspace{3em}foncteur bordant
% archive sheet number: 10
